\section{The complete first-exit surface source}
\label{sec:v83-surface}
We now estimate a quantity defined on the actual outgoing surface,
not a bounded multiplier substituted for that surface. The estimate
uses both a complete starting-source bound and a two-sided payment
of optical variation. It is local in roof mass; the distinction from
pointwise height will be kept explicit.

Let $\tau(x)$ be the forward lifetime in the original regular exact
event and put
\begin{equation}\label{eq:v83-unweighted-exit}
 e_h(t)=\int\one_{\{\tau(x)\le h\}}\,d\mu_t(x).
\end{equation}
This is the \emph{unweighted} starting mass; it keeps both section
tests and the exact label. There is no renormalization after removal
of a boundary layer.

\begin{theorem}[All true exits in a roof interval]
\label{thm:v83-exit-local-mass}
There is $r_0>0$ such that, for every fixed auxiliary band $B_*$,
all $0<r<r_0$, $m\ge2$, and intervals $J$,
\begin{align}
 m^2\int_J e_r(t)\,dt
 \le{}& C r^{1/32}(|J|+B_*^{-1})\notag\\
 &+C_{B_*}(1+B_*|J|)m^3\rho_{B_*}^{m/2}.
                                      \label{eq:v83-exit-finite}
\end{align}
The constants are uniform in $R,n,k$. Consequently, for fixed $r,L$,
\begin{equation}\label{eq:v83-exit-ordered-mass}
 \limsup_m\sup_{R,n,k}\sup_{|J|=L}m^2\int_J e_r(t)\,dt
                                      \le C Lr^{1/32}.
\end{equation}
The same conclusion holds for the $h$-average of $e_h$.
\end{theorem}
\begin{proof}
The all-contact motion estimate places every finite exit of lifetime
at most $r$ in the union of the starting layers
\[
 c_j\le C_Eq_0^{d_j}\sqrt r,\qquad
 \Delta_j\le C_Eq_0^{d'_j}\sqrt r,\qquad
 \operatorname{dist}(T_R^jx,\partial Y_R^*)
                         \le C_Eq_0^{d_j}\sqrt r.
\]
The union includes every contact and flight, and all initial, internal
and terminal decisions. A positive-margin endpoint would continue
the physical word and field; a regular critical point cannot be a
forward exit. Coordinate seams are not stopping faces.

The first two layer unions are contained in the set on which the
original guard with auxiliary parameter $C\sqrt r$ has value zero:
$q_0^d\le q_0^{d/4}$. Their source is therefore bounded by the
complete physical bad source in Input~\ref{input:v83-local}. For
the section layers use the marked transverse-strip estimate there
at the actual collision time $j$, with width
$C_Eq_0^{d_j}\sqrt r$. Each positive upper comparison drops endpoint
tests only on its upper side and costs exactly $c_*^{-1}$; the
occupation and displacement coefficient remain $(n,k)$.

There are at most two contact marks of each depth. The leading
section terms sum as
$Cr^{1/32}\sum_{d\ge0}q_0^{d/16}$, and the physical guard term
has the same power of $r$. The at most $m+1$ section remainders,
together with the physical remainder, are bounded by the second
line of \eqref{eq:v83-exit-finite}. Thus no uncontrolled word count
or sum of collision-dependent bandwidths occurs. Take the collision
limit with $r,L,B_*$ fixed and then enlarge $B_*$. Finally
$e_h\le e_r$ for $h\le r$.
\end{proof}

\subsection{The physical surface measure and its starting-roof assignment}
Use one-sided nongrazing outgoing faces of the original exact event.
On a transverse angular chart a face is $F=b(\vartheta)$, and its
physical flux is
\begin{equation}\label{eq:v83-surface-flux}
 d\zeta=\frac{\omega(b(\vartheta),\vartheta)}
                 {4\pi c_*\chi(b(\vartheta),\vartheta)}\,d\vartheta.
\end{equation}
A countable disjoint refinement counts each face with its incident
physical component. Tangential strata have zero transverse flux.
Let $\tau^-(y)$ be the backward lifetime in that same one-sided event.
Define
\begin{equation}\label{eq:v83-untruncated-flux}
 \mathcal F_h(t)=\int_{\Gamma_{\rm ng}^+}
  \one_{\{0<L(y)-t\le h,\ L(y)-t<\tau^-(y)\}}\,d\zeta(y).
\end{equation}
There is no initial lower bound on $|g|$. The selection in
\eqref{eq:v83-untruncated-flux} excludes later re-entries. Competing-hit,
section and other true exact-event faces remain present. This is the
increasing limit of the original $\mathcal F_h^a$ as $a\downarrow0$
at fixed count, on the common versions of Appendix~\ref{app:v83-kernels}.
The measure is initially allowed to be extended nonnegative.

\begin{theorem}[Two-sided optical payment for the actual flux]
\label{thm:v83-two-sided-flux}
For every $0<\eta<1/2$, jointly for almost every $(t,h)$,
\begin{align}
 \mathcal F_h(t)&\le C e_h(t)+
                C_\eta\mathsf P_{\eta,h}v^\eta(t),
                                     \label{eq:v83-flux-upper}\\
 E_h^\varepsilon(t)+O_h^\varepsilon(t)&\le C\mathcal F_h(t)+
                C_\eta\mathsf P_{\eta,h}v^\eta(t).
                                     \label{eq:v83-source-upper}
\end{align}
In particular the full outgoing assignment is locally integrable in
the starting roof at fixed count. The first inequality is a bound
for the surface source itself, not just for its pullback cost.
\end{theorem}
\begin{proof}
Along a nongrazing first-exit orbit write $w(s)=\omega(\Phi_sx)$
and let $\tau\le h$ be its lifetime. In the angular parametrization,
the starting fiber density is $w(0)/(4\pi c_*\chi(x))$ and the
exit flux density is $w(\tau-)/(4\pi c_*\chi(\Phi_\tau x))$.
Both limiting factors are positive at a nongrazing exit. The two
inequalities
\[
 w(\tau-)\le w(0)+\int_0^\tau[w'(s)]_+\,ds,\qquad
 w(0)\le w(\tau-)+\int_0^\tau[-w'(s)]_+\,ds
\]
are used in their respective directions. The uniform upper and lower
bounds for $\chi$ compare their fiber normalizations.

For \eqref{eq:v83-flux-upper}, sum the first inequality on the actual
first-exit graph. Its starting term is at most $Ce_h$. Pay its
positive variation by \eqref{eq:v83-fractional-drift} and common
\emph{target-direction} coarea covariance. This yields the potential
term. A boundary point is counted with precisely the initial states
selected in \eqref{eq:v83-untruncated-flux}; no multiplicity from a
later return is introduced.

For \eqref{eq:v83-source-upper}, use the second inequality, with
$W^\varepsilon(x)\le1$, for nongrazing exits. At grazing exits
$w(\tau-)=0$ and negative variation pays the full starting weight.
On survivors it pays the positive optical loss. These classes are
disjoint, so their common target insertion is dominated by the full
$v^\eta$ source, not a word-cardinality multiple of it. Truncate to
compact regular face pieces and flow boxes, and pass to the full
nonnegative sums at fixed count. Appendix~\ref{app:v83-kernels}
gives the joint versions and first-exit identification. Local
integrability follows from $e_h\le p$ and
Theorem~\ref{thm:v83-potential-transport}.

The exact pullback identity still contains the inverse forward
Jacobian. The two variation inequalities account for its cost;
they do not replace that Jacobian by one.
\end{proof}

\begin{theorem}[A complete local-mass estimate for outgoing first exits]
\label{thm:v83-surface-local-mass}
For every previously fixed $0<r<r_0$, $L>0$ and $0<\eta<1/2$,
\begin{equation}\label{eq:v83-surface-local-mass}
 \limsup_m\sup_{R,n,k}\sup_{|J|=L}m^2
       \int_J\overline{\mathcal F}_r(t)\,dt
 \le C Lr^{1/32}+C_\eta(L+r)r^{(1+\eta)/2}.
\end{equation}
In particular, on roof intervals of length $r$, the right side is
$O_\eta(r^{33/32})$. Every initial momentum and every nongrazing
one-sided first-exit component is included. This is an estimate of
local surface-assignment mass, not an essential-height estimate.
\end{theorem}
\begin{proof}
Average \eqref{eq:v83-flux-upper} and use $e_h\le e_r$ and
$\mathsf P_{\eta,h}v^\eta\le\mathsf P_{\eta,r}v^\eta$. Fubini gives
\[
 \int_J\mathsf P_{\eta,r}v^\eta(t)\,dt
 \le\frac{r^{1-\kappa_\eta}}{1-\kappa_\eta}
                                  \int_{J+(0,r)}v^\eta(y)\,dy.
\]
Use \eqref{eq:v83-exit-ordered-mass} for the first term and the
fixed-width bound for the actual $v^\eta$ source for the second.
All intervals have fixed lengths before the collision limit.
Since $1-\kappa_\eta=(1+\eta)/2>1/32$, setting $L=r<1$ proves
the last assertion. The surface estimate has been proved through
its first-exit graph and positive optical variation; no theorem
for a bounded volume multiplier is being relabeled as a trace theorem.
\end{proof}
